\renewcommand{\thesubsection}{}

\hypertarget{introduction}{%
\subsection*{Introduction}\label{introduction}}

The Newick Standard for representing trees in computer-readable form
makes use of the correspondence between trees and nested parentheses,
noticed in 1857 by the famous English mathematician
Arthur
Cayley. If we have this rooted tree:

\begin{figure}[!hbt]
	\begin{center}
		\includegraphics[scale=0.4]{img/1876}
		\label{fig:nw}
	\end{center}
\end{figure}

then in the tree file it is represented by the following sequence of
printable characters:

\bigskip

\texttt{(B,(A,C,E),D);}

\bigskip

The tree ends with a semicolon. The bottommost node in this tree is an
interior node, not a tip. Interior nodes are represented by a pair of
matched parentheses. Between them are representations of the nodes that
are immediately descended from that node, separated by commas. In the
above tree, the immediate descendants are B, another interior node, and
D. The other interior node is represented by a pair of parentheses,
enclosing representations of its immediate descendants, A, C, and E. In
our example these happen to be tips, but in general they could also be
interior nodes and the result would be further nestings of parentheses,
to any level.

\bigskip

Tips are represented by their names. A name can be any string of
printable characters except blanks, colons, semicolons, parentheses, and
square brackets.

\bigskip

Because you may want to include a blank in a name, it is assumed that an
underscore character ("\_") stands for a blank; any of these in a name
will be converted to a blank when it is read in. Any name may also be
empty: a tree like

\bigskip

\texttt{(,(,,),);}

\bigskip

is allowed. Trees can be multifurcating at any level.

\bigskip

Branch lengths can be incorporated into a tree by putting a real number,
with or without decimal point, after a node and preceded by a colon.
This represents the length of the branch immediately below that node.
Thus the above tree might have lengths represented as:

\bigskip

\texttt{(B:6.0,(A:5.0,C:3.0,E:4.0):5.0,D:11.0);}

\bigskip

The tree starts on the first line of the file, and can continue to
subsequent lines. It is best to proceed to a new line, if at all,
immediately after a comma. Blanks can be inserted at any point except in
the middle of a species name or a branch length.

\bigskip

The above description is actually of a subset of the Newick Standard.
For example, interior nodes can have names in that standard. These names
follow the right parenthesis for that interior node, as in this example:

\bigskip

\texttt{(B:6.0,(A:5.0,C:3.0,E:4.0)Ancestor1:5.0,D:11.0);}

\hypertarget{examples}{%
\subsection*{Examples}\label{examples}}

To help you understand this tree representation, here are some trees in
the above form:

\bigskip

\texttt{((raccoon:19.19959,bear:6.80041):0.84600,((sea\_lion:11.99700,\ \\ seal:12.00300):7.52973,((monkey:100.85930,cat:47.14069):20.59201,\ \\ weasel:18.87953):2.09460):3.87382,dog:25.46154);}
 
\bigskip

\noindent \texttt{(Bovine:0.69395,(Gibbon:0.36079,(Orang:0.33636,(Gorilla:0.17147,\\ (Chimp:0.19268,\ Human:0.11927):0.08386):0.06124):0.15057):0.54939,\\ Mouse:1.21460):0.10;}

\bigskip

\noindent
\begin{adjustbox}{center}
\texttt{(Bovine:0.69395,(Hylobates:0.36079,(Pongo:0.33636,(G.\_Gorilla:0.17147,}
\end{adjustbox}
\begin{adjustbox}{center}
\texttt{(P.\_paniscus:0.19268,H.\_sapiens:0.11927):0.08386):0.06124):0.15057):0.54939,}
\end{adjustbox}
\noindent \texttt{Rodent:1.21460);}

\bigskip

\texttt{A;}

\bigskip

\texttt{((A,B),(C,D));}

\bigskip

\texttt{(Alpha,Beta,Gamma,Delta,,Epsilon,,,);}

\hypertarget{non-uniqueness}{%
\subsection*{(Non-)Uniqueness}\label{non-uniqueness}}

The Newick Standard does not make a unique representation of a tree, for
two reasons. First, the left-right order of descendants of a node
affects the representation, even though it is biologically
uninteresting. Thus, to a biologist

\bigskip

\texttt{(A,(B,C),D);}

\bigskip

is the same tree as

\bigskip

\texttt{(A,(C,B),D);}

\bigskip

which is in turn the same tree as

\bigskip

\texttt{(D,(C,B),A);}

\bigskip

and that is the same tree as

\bigskip

\texttt{(D,A,(C,B));}

\bigskip

and

\bigskip

\texttt{((C,B),A,D);}

\hypertarget{rooted-and-unrooted-trees}{%
\subsection*{Rooted and unrooted trees}\label{rooted-and-unrooted-trees}}

In addition, the standard is representing a rooted tree. For many
biological purposes we may not be able to infer the position of the
root. We would like to have a representation of an unrooted tree when
decribing inferences in such cases. Here the convention is simply to
arbitrarily root the tree and report the resulting rooted tree. Thus

\bigskip

\texttt{(B,(A,D),C);}

\bigskip

would be the same unrooted tree as

\bigskip

\texttt{(A,(B,C),D);}

\bigskip

and as

\bigskip

\texttt{((A,D),(C,B));}

\hypertarget{widespread-use}{%
\subsection*{Widespread use}\label{widespread-use}}

In spite of this limitation of nonuniqueness the readability of the
resulting representation (for trees of modest size) and the ease of
writing programs that read it have kept this standard in widespread use.

\bigskip

Its competitors include the
NEXUS standard for trees
(part of the more general NEXUS standard for phylogeny data sets).
However the NEXUS representation of trees is based on the Newick
standard -\/- inside the NEXUS TREES Block you will find ... Newick
trees.

\bigskip

A less Newick-based standard is the
PhyloXML standard, which
is an XML representation using nesting the \textless CLADE\textgreater{}
... \textless/CLADE\textgreater{} tag pairs instead of parentheses.

\hypertarget{origin}{%
\subsection*{Origin}\label{origin}}

The Newick Standard was adopted 26 June 1986 by an informal committee
meeting convened by me during the
Society for the Study of
Evolution meetings in Durham, New Hampshire and consisting of
James Archie,
William H.E.
Day,
Wayne
Maddison, Christopher
Meacham, F. James Rohlf,
David
Swofford, and
myself. (The
committee was not an activity of the SSE nor endorsed by it). The reason
for the name is that the second and final session of the committee met
at Newick's restaurant in Dover, New
Hampshire, and we enjoyed the meal of lobsters. The tree representation
was a generalization of one developed by Christopher Meacham in 1984 for
the tree plotting programs that he wrote for the
PHYLIP package
while visiting Seattle. His visit was a sabbatical leave from the
University of Georgia, which thus indirectly partly funded that work.

%\hypertarget{other-descriptions-of-the-newick-standard}{%
%\subsection*{Other descriptions of the Newick
%Standard}\label{other-descriptions-of-the-newick-standard}}
%
%There has been no formal publication of the Newick Standard.
%
%\begin{itemize}
%\tightlist
%\item
%  {Gary Olsen} has produced a formal description of it which is
%  available \href{newick_doc.html}{here}.
%\item
%  There is a Wikipedia page on the Newick Standard available
%  \href{http://en.wikipedia.org/wiki/Newick_format}{here}
%\end{itemize}
%
\begin{center}\rule{0.5\linewidth}{0.5pt}\end{center}

%\href{../phylip.html}{\includegraphics{icons/PHYLIP.gif} ... to the PHYLIP home page}

Source:  Joseph Felsenstein and Gary Olsen at

\href{http://evolution.genetics.washington.edu/phylip/newicktree.html}{\texttt{\small http://evolution.genetics.washington.edu/phylip/newicktree.html}}